% ===== 原文 pp.119--124 =====
% [原文 p.119]
\begin{figure}[htbp]
\centering
\includegraphics[width=0.96\textwidth]{fig27.png}
\figtitle{27}{生成 0-形式对称性 $g\in G^{(0)}$ 的余维-1 拓扑算符，作用于生成 1-形式对称性 $\gamma\in G^{(1)}$ 的余维-2 拓扑算符，并把它转变成生成 1-形式对称性 $\rho(g)\cdot\gamma\in G^{(1)}$ 的余维-2 拓扑算符。}
\end{figure}

该 Postnikov 类刻画 1-形式对称性的背景场 $B_2$ 的不闭合性
\begin{equation}
  \delta B_2+B_1^*\Theta=0\,, \tag{5.21}
\end{equation}
其中 $B_1$ 是 0-形式对称性的背景场。

\para{更一般的 2-群。}在上面的讨论中，我们没有纳入 0-形式对称性对 $p$-形式对称性的非平凡作用。见图 27。把这类作用纳入 2-群情形后，我们发现最一般的离散 2-群 $\mathfrak{G}^{(2)}$ 可指定为
\begin{equation}
  \mathfrak{G}^{(2)}=\big\{G^{(0)},G^{(1)},\rho,[\Theta]\big\}\,, \tag{5.22}
\end{equation}
其中 $G^{(0)}$ 与 $G^{(1)}$ 是离散的 0-形式与 1-形式对称群，$\rho$ 是 $G^{(0)}$ 在 $G^{(1)}$ 上的作用，它是一个同态
\begin{equation}
  \rho:\ G^{(0)}\to\operatorname{Aut}\left(G^{(1)}\right), \tag{5.23}
\end{equation}
最后 \textbf{Postnikov 类}是
\begin{equation}
  [\Theta]\in H^3_\rho\left(G^{(0)},G^{(1)}\right), \tag{5.24}
\end{equation}
其中该上同调群被 $\rho$ 扭化。

% [Tutor Note: (5.22) 中的 $\mathfrak{G}^{(2)}$ 用 Fraktur 体，以区别于普通对称群 $G$；(5.24) 的下标 $\rho$ 表示系数被 $\rho$ 扭化，即用下面 (5.26) 的微分 $\delta_\rho$ 代替普通上链微分 $\delta$。校验：当 $\rho$ 平凡时 $H^3_\rho$ 退化为普通 $H^3(G^{(0)},G^{(1)})$，与 (5.20) 一致。]

% [原文 p.120]
\begin{notebox}{扭化群上同调}
给定形如
\begin{equation}
  \rho:\ G^{(0)}\to\operatorname{Aut}(G^{(1)})\,, \tag{5.25}
\end{equation}
的作用，我们可以把群上同调的微分修改为
\begin{align}
  \delta_\rho\alpha_p(g_1,g_2,\cdots,g_{p+1})
  &=\rho(g_1)\cdot\alpha_p(g_2,g_3,\cdots,g_{p+1})\alpha_p^{s(p+1)}(g_1,g_2,\cdots,g_p)\nonumber\\
  &\quad\ \prod_{i=1}^{p}\alpha_p^{s(i)}(g_1,g_2,\cdots,g_{i-1},g_ig_{i+1},g_{i+2},g_{i+3}\cdots,g_{p+1})\,, \tag{5.26}
\end{align}
其中 $\alpha_p$ 是 $G^{(0)}$ 上取值于 $G^{(1)}$ 的 $p$-群上链，而
\begin{equation}
  \rho(g_1)\cdot\alpha_p(g_2,g_3,\cdots,g_{p+1})\,, \tag{5.27}
\end{equation}
表示自同构 $\rho(g_1)\in\operatorname{Aut}(G^{(1)})$ 对 $\alpha_p(g_2,g_3,\cdots,g_{p+1})\in G^{(1)}$ 的作用。

这样得到的上同调群记作
\begin{equation}
  H^p_\rho\left(G^{(0)},G^{(1)}\right), \tag{5.28}
\end{equation}
并被称为 $G^{(0)}$ 取值于 $G^{(1)}$、扭化为 $\rho$ 的第 $p$ 个\textbf{扭化群上同调}。
\end{notebox}

% [Tutor Note: (5.26) 符号 $s(i)$ 已在 (4.46)（p.70）定义为 $s(n):=(-1)^n$，故这里不必另行推断。据此 (5.26) 就是扭化群上链的显式公式：首项为 $i=0$ 的「$+$」项、次项为 $i=p+1$ 的 $(-1)^{p+1}$ 项、求和项逐个带上 $(-1)^i$，与 $\delta\alpha=\sum_i(-1)^i\alpha(\cdots)$ 的标准形式一致。核对：取 $\rho=1$ 时 (5.26) 退化为普通群上链微分。]

\para{分裂 2-群。}当作用 $\rho$ 非平凡但 Postnikov 类消失
\begin{equation}
  [\Theta]=0\in H^3_\rho\left(G^{(0)},G^{(1)}\right), \tag{5.29}
\end{equation}
时，所得的 2-群称为\textbf{分裂 2-群}。

\begin{examplebox}{例 5.1：纯 $\SU(N)$ Yang--Mills 中的电荷共轭}
考虑 $d$ 维纯 $\SU(N)$ Yang--Mills 理论，它具有电 1-形式对称性
\begin{equation}
  G^{(1)}=\Z_N\,. \tag{5.30}
\end{equation}
该理论还含有一个
\begin{equation}
  G^{(0)}=\Z_2\,, \tag{5.31}
\end{equation}
电荷共轭 0-形式对称性，它对应于 $\SU(N)$ 的、翻转 Dynkin 图的外自同构。

这两个对称性组合成一个分裂 2-群，因为 $\SU(N)$ 的外自同构作为其 $\Z_N$ 中心的外自同构作用，使得
\begin{equation}
  p\longmapsto-p\,,\qquad p\in\{0,1,2,\cdots,N-1\}\,. \tag{5.32}
\end{equation}
\end{examplebox}

% [Tutor Note: (5.32) 中的映射 $p\mapsto -p$ 定义在 $\Z_N$ 的整数代表元集合 $\{0,1,\cdots,N-1\}$ 上，故 $-p$ 应按模 $N$ 理解。电荷共轭把 $\SU(N)$ 的 $N$ 维表示变成其复共轭表示，从而把中心荷 $p$ 变为 $-p$，这就是「$\rho$ 非平凡、但 $[\Theta]=0$」的分裂 2-群的实例。]

% [原文 p.121]
\subsubsection{与混合 't Hooft 反常的关系}

考虑一个高阶群，其中唯一的不闭合关系是 (5.19)，并假定该高阶群对称性的 't Hooft 反常使得它在组成的 $G^{(p)}$ $p$-形式对称群上平凡化。在这种情形下，我们可以规范化 $G^{(p)}$。在规范化后的理论中，我们得到对偶的
\begin{equation}
  G^{(d-p-2)}=\widehat{G}^{(p)}\,, \tag{5.33}
\end{equation}
$(d-p-2)$-形式对称性。$G^{(p)}$ 参与高阶群对称性 (5.19) 这一事实，被替换为对偶的 $G^{(d-p-2)}$ 对称性参与一个混合 't Hooft 反常这一事实，相应的反常理论为
\begin{equation}
  \cI_{d+1}=\exp\left(-2\pi\ii s(d-p)\int B_{d-p-1}\cup_\eta\Theta(B_{p_1+1},B_{p_2+1},\cdots B_{p_n+1})\right),\qquad
  s(d-p)=(-1)^{d-p}\,. \tag{5.34}
\end{equation}
这段论证与把 (4.194) 变成 (4.197) 的那段论证相同，后者已在群扩张的语境中讨论过。

% [Tutor Note: (5.34) 是群扩张结论 (4.194)→(4.197) 的高阶形式版本：规范化 $G^{(p)}$ 后，原高阶群的「不闭合」信息 $(5.19)$ 变成对偶 $(d-p-2)$-形式对称性的混合 't Hooft 反常。符号 $s(d-p)$ 的来历：$(d-p)$ 维中交换两个算符需要 $(-1)^{d-p}$ 的交换因子，故反常的符号随维数奇偶而变。]

\subsection{混合连续-离散 2-群对称性}

上文中我们讨论了所有参与的高阶形式对称性都连续的高阶群，以及都离散的高阶群。也可以有一些高阶形式对称性连续、另一些高阶形式对称性离散的高阶群。

在本小节中，我们讨论 0-形式对称群 $G^{(0)}$ 连续而 1-形式对称群 $G^{(1)}$ 离散的 2-群对称性。这类 2-群的定义关系是
\begin{equation}
  \delta B_2+B_1^*\Theta=0\,, \tag{5.35}
\end{equation}
其中 $\Theta$ 是一个类
\begin{equation}
  [\Theta]\in H^3(BG^{(0)},G^{(1)})\,, \tag{5.36}
\end{equation}
的代表，这个类再次被称为 \textbf{Postnikov 类}，其中 $H^3(BG^{(0)},G^{(1)})$ 是 $G^{(0)}$ 的分类空间 $BG^{(0)}$ 的、以 $G^{(1)}$ 为系数的三阶上同调群，而 $B_1^*\Theta$ 是 $\Theta$ 在映射 $B_1$（从时空到 $BG^{(0)}$，描述 0-形式背景的主丛）下的拉回。

% [Tutor Note: (5.35) 与离散情形的 (5.21) 形式完全相同，(5.21) 中 $B_1$ 是 0-形式对称群 $G^{(0)}$ 的背景场；$K^{(0)}$ 的背景场 $B_1^K$ 只出现在 §5.2.1 的 (5.18)。此处 $B_1$ 是连续群 $G^{(0)}$ 主丛的分类映射 $M\to BG^{(0)}$；因此 $\Theta$ 的取值域从群上同调 $H^3(G^{(0)},G^{(1)})$ 变为分类空间上同调 $H^3(BG^{(0)},G^{(1)})$（对照 (5.20) 与 (5.36)）。]

\begin{notebox}{分类空间}
给定某个群 $G$ 的分类空间 $BG$，可以从具有如下两条性质的空间 $EG$ 出发来得到：
\begin{itemize}[leftmargin=1.4em,itemsep=1pt,topsep=1pt]
\item $EG$ 的所有同伦群都平凡；
% [原文 p.122]
\item $G$ 在 $EG$ 上有自由作用（即没有不动点）。
\end{itemize}

然后我们令
\begin{equation}
  BG:=EG/G\,, \tag{5.37}
\end{equation}
其中对 $EG$ 取商是由该自由 $G$ 作用实现的。

注意分类空间 $BG$ 并不唯一。给定固定的群 $G$，可以有多个满足上述判据的空间 $EG$ 与 $BG$。

注意 $EG$ 自然地是 $BG$ 上的主 $G$-丛，并被称为万有 $G$-丛，因为可以把任意流形 $M$ 上的任意主丛 $P$ 作为拉回
\begin{equation}
  P=f^*(EG)\,, \tag{5.38}
\end{equation}
得到，即该万有丛 $EG$（$BG$ 上的）在映射
\begin{equation}
  f:\ M\to BG \tag{5.39}
\end{equation}
下的拉回。
\end{notebox}

% [Tutor Note: 分类空间只定义到弱等价，故「不唯一」；$EG\to BG$ 是万有主丛，$P=f^*(EG)$ 说明主丛由映射 $f$ 的同伦类分类，即 $\operatorname{Prin}_G(M)\cong[M,BG]$。核对 (5.37)：$EG$ 可缩且 $G$ 自由作用，故 $\pi_n(BG)=\pi_{n-1}(G)$，$BG$ 的同伦型唯一。]

我们将专注于研究 \textbf{Bockstein 2-群}，它们在各种维度的规范理论中非常自然地出现 [186,188]。

\subsubsection{0-形式对称性的全局形式}

在描述 Bockstein 2-群之前，我们必须先说明关于连续李群 0-形式对称性的精确全局形式的一些预备知识。我们将采用 [207] 第 2 节所讨论的途径。

\para{0-形式对称代数。}考虑基于一个紧致半单非阿贝尔李代数的连续 0-形式对称性
\begin{equation}
  \mathfrak{f}=\text{0-形式对称性的李代数}\,. \tag{5.40}
\end{equation}
这个 0-形式代数正是 Noether 流局域算符
\begin{equation}
  j^\mu_a(x)\,, \tag{5.41}
\end{equation}
的代数，其中 $a$ 是 $\mathfrak{f}$ 中的指标，而 $\mu$ 是时空指标。

\para{0-形式对称群。}令
\begin{equation}
  F=\text{与 }\mathfrak{f}\text{ 相关联的单连通李群}\,. \tag{5.42}
\end{equation}
可能并非 $F$ 的所有元素都在理论上非平凡地作用，这意味着 $F$ 在 4.2.2 节所讨论的意义上不是完全忠实的。一般来说，一个子群
\begin{equation}
  \mathcal{Z}\subseteq Z(F)\,, \tag{5.43}
\end{equation}
% [原文 p.123]
可能是非忠实的，其中 $Z(F)$ 是 $F$ 的中心。于是我们有以下定义。

\begin{definitionbox}{定义 5.1：0-形式对称群}
完全忠实作用的 0-形式对称群是
\begin{equation}
  G^{(0)}=F/\mathcal{Z}\,, \tag{5.44}
\end{equation}
它被简称为理论的 \textbf{0-形式对称群}，有时也称为连续 0-形式对称性的全局形式。
\end{definitionbox}

我们如何确定群 $\mathcal{Z}$？这可以通过研究 $\mathfrak{f}$ 对真局域算符的作用来完成，这些算符构成 $\mathfrak{f}$ 的表示 $R_i$。于是 $\mathcal{Z}$ 是
\begin{equation}
  \mathcal{Z}=\text{保持所有局域算符表示 }R_i\text{ 不变的 }Z(F)\text{ 的子群}\,. \tag{5.45}
\end{equation}
把离散 0-形式对称性与连续 0-形式对称性一并纳入，会导向不连通的 0-形式对称群。详见 [198]。

% [Tutor Note: (5.42)--(5.45) 的记号要分清：$F$ 是单连通群、$Z(F)$ 是其中心（普通体字 $Z$），而 $\mathcal{Z}$ 是其中不忠实作用的子群（花体）。校验 (5.44)：$G^{(0)}=F/\mathcal{Z}$ 必与 $F$ 有同一李代数 $\mathfrak{f}$，且是忠实作用的最大商群；例 5.2 中 $\mathcal{Z}=Z(F)=\Z_2$ 给出 $\SU(2)/\Z_2=\SO(3)$。]

\para{0-形式对称丛。}连续 0-形式对称性的背景场是 0-形式对称群 $G^{(0)}$ 的主丛上的联络。可能的丛对全局形式很敏感。特别地，一个 $G^{(0)}$ 丛可能不可提升为一个 $F$ 丛，其障碍类为 $\left[w_2^f\right]$，取值于 $\mathcal{Z}$。

\para{规范理论中的味对称群。}在规范理论中，局域算符来自场强与物质场的组合。如我们先前所讨论的，若这样的组合在规范群 $\mathcal{G}$ 下变换，则它给出一个生活在 Wilson 线算符端点处的非真局域算符。因此，真局域算符对应于场强与物质场的规范不变组合。这引向如下计算。

设 $\phi_i$ 是物质场，它们在 $\mathcal{G}$ 的不可约表示 $R_i^g$ 与 $F$ 的不可约表示 $R_i^f$ 中变换。这一信息给出荷
\begin{equation}
  \left(q_i^g,q_i^f\right)\in\widehat{Z}(\mathcal{G})\times\widehat{Z}(F)\,. \tag{5.46}
\end{equation}
这些荷生成一个子群
\begin{equation}
  Q\subseteq\widehat{Z}(\mathcal{G})\times\widehat{Z}(F)\,. \tag{5.47}
\end{equation}

% [Tutor Note: (5.46) 的物理含义：$R_i^g$ 是规范群 $\mathcal{G}$ 的表示、$R_i^f$ 是味群 $F$ 的表示，而 $\widehat{Z}(\mathcal{G})$ 与 $\widehat{Z}(F)$ 是相应中心的特征标群（Pontryagin 对偶），故荷对 $(q_i^g,q_i^f)$ 记录的是「表示在中心元下取到的相位」。校验：若表示 $R_i$ 在中心元 $z$ 下给出相位 $\chi_i(z)$，则 $\chi_i\in\widehat{Z}$；多个物质场在同一表示下的荷相加，故物质场张成的荷集合构成子群 $Q$。]

% [原文 p.124]
规范不变算符的 0-形式中心荷是
\begin{equation}
  Q_{\mathrm{genuine}}=Q\cap\widehat{Z}(F)\subseteq\widehat{Z}(F)\,. \tag{5.48}
\end{equation}
我们所要的群 $\mathcal{Z}\subseteq Z(F)$ 则是与 $Q_{\mathrm{genuine}}$ 平凡配对的子群，其 Pontryagin 对偶可计算为
\begin{equation}
  \widehat{\mathcal{Z}}=\frac{\widehat{Z}(F)}{Q_{\mathrm{genuine}}}\,. \tag{5.49}
\end{equation}

% [Tutor Note: (5.48)--(5.49) 的逻辑链：物质场荷生成 $Q\subseteq\widehat{Z}(\mathcal{G})\times\widehat{Z}(F)$；真算符（规范不变）的中心荷只是 $Q$ 中落在 $\widehat{Z}(F)$ 的部分，即 $Q_{\mathrm{genuine}}$；与这些荷平凡配对的中心子群就是实际忠实作用的 $\mathcal{Z}$，故 $\widehat{\mathcal{Z}}=\widehat{Z}(F)/Q_{\mathrm{genuine}}$。用下面的例 5.2 可数值校验：$Q=\langle(2,1)\rangle\subset\Z\times\Z_2$，与 $\widehat{Z}(F)=\Z_2$ 的交为 $0$，故 $\widehat{\mathcal{Z}}=\Z_2$、$\mathcal{Z}=\Z_2$。]

\begin{examplebox}{例 5.2：阿贝尔规范理论}
考虑一个 $d$ 维的
\begin{equation}
  \mathcal{G}=U(1)\,, \tag{5.50}
\end{equation}
规范理论，它带有 2 个规范荷为
\begin{equation}
  q^g=2 \tag{5.51}
\end{equation}
的无质量复标量。存在一个
\begin{equation}
  \mathfrak{f}=\mathfrak{su}(2)\,, \tag{5.52}
\end{equation}
味对称性代数，它在 $\mathfrak{su}(2)$ 的基础表示中转动这两个复标量。相关联的单连通群是
\begin{equation}
  F=SU(2)\,. \tag{5.53}
\end{equation}
我们有
\begin{equation}
  \widehat{Z}(\mathcal{G})\times\widehat{Z}(F)=\Z\times\Z_2\,, \tag{5.54}
\end{equation}
其中物质场生成子群
\begin{equation}
  Q=\langle(2,1)\rangle\in\Z\times\Z_2\,, \tag{5.55}
\end{equation}
由此我们算出
\begin{equation}
  Q_{\mathrm{genuine}}=Q\cap\widehat{Z}(F)=0\,, \tag{5.56}
\end{equation}
这意味着
\begin{equation}
  \mathcal{Z}=Z(F)=\Z_2\,, \tag{5.57}
\end{equation}
而 0-形式味对称群是
\begin{equation}
  G^{(0)}=F/\mathcal{Z}=SU(2)/\Z_2=SO(3)\,. \tag{5.58}
\end{equation}
\end{examplebox}

% [Tutor Note: 例 5.2 的数值自洽性：(5.51) 中两个复标量的 $\U(1)$ 规范荷为 2，故其在 $\widehat{Z}(\mathcal{G})=\Z$ 上的荷为 2；(5.53) 中 $F=\SU(2)$ 的基础表示在该 $\Z_2$ 上的荷为 1，于是物质场荷生成 $(2,1)$ 及其整数倍，即 (5.55)。校验 (5.57)：$\mathcal{Z}=Z(F)$ 意味着 $\SU(2)$ 的中心 $\Z_2$ 整体不忠实作用，故有效 0-形式群为 $\SO(3)$——这与「两个 $\U(1)$ 电荷为 2 的标量不足以探测 $\SU(2)$ 中心的 $\Z_2$」的物理一致。]

\begin{examplebox}{例 5.3：非阿贝尔规范理论}
考虑一个 $d$ 维非阿贝尔规范理论，其规范群为
\begin{equation}
  \mathcal{G}=\Spin(4N+2)\,, \tag{5.59}
\end{equation}
% 注：本 examplebox 在原文中跨 p.124--p.125，按 Lead 裁决由 tr-08 在其片段开头续写正文并在 (5.69) 之后闭合；
% 因此 tr-07 的片段到此为止，不写该方框的结束标记，以便与 tr-08 的内容合成同一个方框。
